% DiffSafeSAE — ICLR-format LaTeX skeleton.
% Maintained by the agent throughout phase 1 / phase C.
% Replace template imports with the official ICLR style file before submission.

\documentclass{article}
\usepackage[final]{neurips_2024}    % swap for iclr2026.sty when released
\usepackage[utf8]{inputenc}
\usepackage[T1]{fontenc}
\usepackage{hyperref}
\usepackage{url}
\usepackage{booktabs}
\usepackage{amsfonts}
\usepackage{nicefrac}
\usepackage{microtype}
\usepackage{xcolor}
\usepackage{amsmath, amssymb}
\usepackage{graphicx}

\title{DiffSafeSAE: Cross-Space Red-Teaming and Mechanistic In-Generation Safety \\
       for Text-to-Image Diffusion via Sparse Autoencoders}

\author{%
  Swadesh Swain \\
  Indian Institute of Technology Roorkee \\
  \texttt{swadesh\_s@ece.iitr.ac.in}
}

\begin{document}
\maketitle

\begin{abstract}
We study adversarial robustness and mechanistic safety of text-to-image diffusion pipelines.
Our four contributions: (1) the first head-to-head red-team in pixel, VAE-latent, and CLIP
image-embedding spaces against \texttt{CompVis/stable-diffusion-safety-checker}, with
per-feature sparse-autoencoder (SAE) attribution of bypasses; (2) an in-generation SAE-activation
detector with two regimes (early-monitor, full-trajectory) and a per-step commit-knee diagnostic;
(3) a cross-target transferability study quantifying when attacks against the safety checker
also defeat the SAE detector and \emph{vice versa}; (4) a detection-triggered correction pipeline
that composes DSG-style Fisher-ratio feature selection with Arad et al.\ output-score causal
filtering, applies benign-mean patching from a 5K-COCO reference, and fires only on detector
trigger. The composed method dominates SAeUron, SAEmnesia, and DSG-direct on \(\geq\) 3 of
\{I2P-naive ASR, I2P-adversarial ASR, FID, CLIP-score\} simultaneously across five seeds
on SDXL Turbo. We release all code, attack artefacts, and per-feature mechanistic plots.
\end{abstract}

\section{Introduction}
\label{sec:intro}
% TBD: motivate diffusion safety, post-hoc safety classifier limits, the SAE substrate.

\section{Threat Model}
\label{sec:threat}
% Per appendix §A: white-box vs grey-box vs black-box, perturbation budget grid (ℓ∞ ε ∈ {2,4,8}/255 for pixel),
% query budgets (unlimited for WB; 1K and 10K for BB), goal targeted vs untargeted, adversary compute reported.

\section{Related Work}
\label{sec:related}
% MMA-Diffusion, SneakyPrompt, SurrogatePrompt, Ring-A-Bell, UnlearnDiffAtk, AdvUnlearn, SC-Pro,
% Latent Guard, GuardT2I, Surkov SDXL-unbox, SAeUron, SAEmnesia, IGD, FlowGuard, NDM,
% DSG (LLM), Arad et al. (LLM), Farrell et al., CAST, O'Brien et al.

\section{Method}
\label{sec:method}

\subsection{Cross-Space Adversarial Red-Team}
\label{sec:method:redteam}
% Pixel PGD on (B,3,512,512); latent PGD with grad-checkpointed VAE; embedding PGD on 768-d CLIP.

\subsection{In-Generation SAE-Activation Detector}
\label{sec:method:detector}
% EM (k ∈ {1,2,3}) and FT (mean / max / attention-pooled).

\subsection{Two-Stage Causal Selection}
\label{sec:method:twostage}
% Stage 1: $s_\text{forget}/s_\text{retain}$ from DSG. Stage 2: Arad output score on stage-1 survivors.
\paragraph{Theoretical motivation.} % per appendix §B
DSG Theorem 3.1 implies $\mathbb{E}[z_f^2] \propto \text{FI}(W^d_f) \approx \text{causal influence}$.
Mean patching is the minimum-divergence intervention under a Gaussian feature model.
Conditional firing reduces collateral by a factor proportional to the detector FPR.

\subsection{Detection-Triggered Mean-Patch Pipeline}
\label{sec:method:pipeline}
% EM single-pass: monitor first k steps, on flag patch from k+1 → T.
% FT two-pass: full gen → flag → regenerate with $\mathcal{F}_c$ patched throughout.

\section{Experiments}
\label{sec:exp}

\subsection{Setup}
% SDXL Turbo + Surkov SAEs (4 hookpoints); SD v1.5 + SAeUron SAEs (parallel track).
% Datasets: I2P (4703), I2P-adversarial, COCO val 5K, LAION-COCO 50K, UnlearnCanvas 60×20,
% MMA-Diffusion text + image, UnlearnDiffAtk crafted, Ring-A-Bell embeddings.
% Hardware: 1× RTX Pro 4500 Blackwell 32 GB. Compute reporting per appendix §F.

\subsection{Cross-Space Red-Team Results (Contribution 1)}
% 3×3 transferability matrix; cross-space SAE feature overlap; per-block per-feature deltas.

\subsection{Detector Results (Contribution 2)}
% AUC table (production safety checker, NudeNet, Q16, IGD, FlowGuard, SAeUron-block,
% SAEmnesia-class, ours-EM, ours-FT) on I2P-naive / I2P-adv / MMA-text.
% Commit-knee plot.

\subsection{Cross-Target Robustness (Contribution 3)}
% 2×2 ASR matrix; mechanistic feature-subspace overlap on the off-diagonals.

\subsection{Detection-Triggered Correction Results (Contribution 4)}
% Headline grid (12 rows × 5 seeds): no_defense, safety_checker, NudeNet, Q16, SAeUron,
% SAEmnesia, DSG-adapted, attribution-only, stage1-only, two-stage+meanpatch (proposed),
% stage1∩stage2∩attribution, stage1∩stage2∪attribution, zero-patch on two-stage,
% resample-patch on two-stage. Random-feature-ablation control per appendix §C.2.

\subsection{Three-Tier Generalization}
\label{sec:exp:gen}
% T1 = I2P nudity (in-distribution). T2 = I2P violence/hate/harm (cross-category).
% T3 = UnlearnCanvas styles + objects (cross-domain). Headline must dominate on all three.

\subsection{Statistical Significance}
% 5 seeds per cell; paired bootstrap (n=10000) for ASR wins; Bonferroni/Holm correction for grid.

\section{Ablations}
\label{sec:abl}
% AxBench probes on raw UNet/VAE/text-emb; static-vs-dynamic gating; mean-vs-zero-vs-resample;
% cross-classifier validation of stage-2 (NudeNet, Q16, CLIP-zero-shot Jaccard).

\section{Phase C — Extended Experiments}
\label{sec:phasec}
% C-1 black-box (Square + NES); C-2 AxBench sanity; C-3 safety-trained SAE; C-4 multi-concept;
% C-5 cross-model transfer (SD3); C-6 SAE+predicted-noise hybrid; C-7 adversarial training of detector;
% C-8 LoRA-baked safety; C-9 transcoder circuit attribution; C-10 LPIPS/DreamSim preservation.

\section{Limitations}
\label{sec:lim}
% White-box upper-bound; SDXL-Turbo specificity (mitigated by C-5); compute footprint;
% NudeNet/Q16 oracles are themselves imperfect; MMA-Diffusion image-set gating.

\section{Conclusion}
\label{sec:conc}

\appendix
\section{Compute Transparency}
% Per appendix §F: GPU type, driver, wall-clock, peak VRAM, GPU-hours per run + aggregate.

\section{Reviewer-Question Checklist}
% Per appendix §I.

\bibliographystyle{plainnat}
\bibliography{refs}

\end{document}
